\section{最小缓存驻留调度}
\subsection{生命周期模型与下界}
一个缓冲区在申请后持续占用逻辑缓存，直到全部使用结束并执行释放。调度的作用是改变这些区间的相互重叠：过早申请会延长区间，过晚释放会保留已经无用的数据，而不同计算分支的交错还会改变末次使用位置。因此本问先处理管理节点的可移动性，再搜索操作节点的相对顺序。

问题一只优化节点顺序，不给缓冲区设置物理偏移，也不加入换入换出。缓冲区从申请后到释放前始终计入驻留量。对一个节点置换 $\pi$，定义处理完序列第 $k$ 个节点后的活跃指示量
\begin{equation}
y_b(k;\pi)=\mathbf 1\{\pi(a_b)\le k<\pi(f_b)\},\qquad k=0,\ldots,N.
\end{equation}
于是调度模型为
\begin{align}
\min_{\pi}\quad&H(\pi)=\max_{0\le k\le N}\sum_{b\in B}s_b y_b(k;\pi),\label{eq:peak}\\
\text{s.t.}\quad&\pi(u)<\pi(v),\qquad (u,v)\in E,\label{eq:topo}\\
&\pi(a_b)<\pi(v)<\pi(f_b),\qquad v\in V_o,\ b\in B_v,\label{eq:use}\\
&\pi:V\longrightarrow\{1,\ldots,N\}\text{ 为双射}.\nonumber
\end{align}
式\eqref{eq:use}把使用时必须活跃的语义显式列出；在输入依赖充分描述生命期时，它由原始路径蕴含。申请记正增量、释放记负增量的题面定义，与式\eqref{eq:peak}完全一致。总峰值是在同一位置对各类型求和后取最大值，不能相加各类型在不同位置达到的峰值。

每条操作的全部输入和输出缓冲区必须同时存在。因此，无论如何重排，都有
\begin{equation}
H^*\ge L_H:=\max_{v\in V_o}\sum_{b\in B_v}s_b.\label{eq:hlower}
\end{equation}
此界利用单个操作的同时使用需求。对任一可行序列，在操作 $v$ 发生前，其缓冲区均已申请且尚未释放，因此该时刻的驻留量至少为右端求和；再对全部操作取最大值得到式\eqref{eq:hlower}。当可行候选达到 $L_H$ 时，上下界闭合，峰值最优性随之成立。多个操作之间必须继续保留的对象会使实际下界更高，故后文同时给出可行上界及这一可直接计算的下界。

也可以从执行前沿理解峰值。将某一操作前缀记为 $P$，其中已经开始使用、仍被后续操作需要的缓冲区构成跨前沿集合。它们把已经完成的计算与尚未完成的计算连接起来，生命期收缩后仍须保留。选择下一个就绪操作，会同时改变首次使用对象和末次使用对象；局部评分衡量这一即时变化，完整峰值则取决于所有前沿中最大的驻留总量。

设当前驻留为 $R$，本步全部新增申请总长为 $\widetilde n_v$，操作后实际释放总长为 $d_v^{\rm real}$。全部新增申请包括使用对象和额外显式申请前驱。该步驻留先增至 $R+\widetilde n_v$，随后降至 $R+\widetilde n_v-d_v^{\rm real}$，历史峰值更新为
\begin{equation}
H_{\rm new}=\max\{H_{\rm old},R+\widetilde n_v\}.
\label{eq:steppeak}
\end{equation}
释放发生在操作之后。即使 $\widetilde n_v-d_v^{\rm real}$ 很小，本步也可能先经历较大的申请峰值。这一区别是局部净增量与全局峰值不能直接等同的原因，也是保留完整重放评价的依据。

\subsection{固定操作顺序下的生命周期收缩}
题面规定申请节点为根、释放节点为叶，六图均满足这一结构。固定非管理操作的相对顺序后，一个申请必须早于其所有显式后继以及所有引用该缓冲区的操作。称其中最早出现的节点为首次需求节点。相应地，释放必须晚于全部显式前驱和全部使用者，称其中最晚出现的节点为末次需求节点。这里不能仅以第一次读或最后一次读代替需求节点，否则可能遗漏额外管理依赖。

\textbf{收缩性质。}在上述前提下，将申请移动到首次需求节点之前，将释放移动到末次需求节点之后，不会增加驻留峰值。

证明分两步。首先，申请是根节点，向后移动不会违反其入边；又因为尚未跨过任何显式后继或使用者，其出边和使用约束仍成立。将它与中间节点逐个交换，只会删除被跨越区间内这一缓冲区的正容量贡献。其次，释放是叶节点，向前移动不会违反其出边；停止在全部前驱和使用者之后，则其入边不受影响。逐个交换同样只会减少中间区间的驻留量。其他缓冲区的申请、释放相对关系不变，因此每一步交换的最大累计量均不增加。对所有缓冲区重复此变换，得到同一操作顺序下不劣的规范化序列。

收缩后，每个对象的区间端点贴近需求范围，搜索因而集中到操作相对顺序。共享对象的释放仍以末次使用和全部显式前驱完成为条件。对于申请有入边或释放有出边的扩展模型，可移动端点还需受到这些边的限制；本文六图满足根申请、叶释放条件。

\subsection{九节点算例：从延迟申请到峰值闭合}
构造三个操作 $u\to v\to w$，缓冲区 $a,b,c$ 的大小分别为 4、3、5。操作 $u$ 使用 $a$，$v$ 使用 $a,b$，$w$ 使用 $b,c$。每个对象分别设申请节点 $A_a,A_b,A_c$ 和释放节点 $F_a,F_b,F_c$；申请指向全部使用者，使用者指向对应释放。此例共九个节点，用于检验生命周期模型，不计入六个题面图的实验统计。

若先申请全部对象，再依次执行三个操作，则申请阶段驻留依次为 4、7、12，峰值已经达到 12。即使此后每次操作完成就尽早释放，该峰值也无法撤销。将同一操作顺序收缩后得到
\begin{equation}
A_a,\ u,\ A_b,\ v,\ F_a,\ A_c,\ w,\ F_b,\ F_c.
\end{equation}
表\ref{tab:toy_life}列出每个节点后的状态，使申请与释放的贡献可以逐步复算。
\begin{table}[H]\centering\small
\caption{构造小例的收缩序列与驻留量}\label{tab:toy_life}
\begin{tabular}{ccccccccc}\toprule
$A_a$&$u$&$A_b$&$v$&$F_a$&$A_c$&$w$&$F_b$&$F_c$\\\midrule
4&4&7&7&3&8&8&5&0\\\bottomrule
\end{tabular}\end{table}

操作 $w$ 同时需要大小为 3 和 5 的两个对象，所以 $L_H=8$。表中峰值恰为 8，因而这一构造例的最优值为 8。对其全部合法拓扑序进行递归枚举，共得到 126 个序列，其中 120 个峰值为 12，6 个峰值为 8，与下界及手算一致。枚举过程中每次只从全部前驱已输出的节点中选择，逐节点更新驻留和峰值，输出终止时检查驻留归零。

这一小例把两个层次区分开：固定操作链时，延迟申请便能排除不必要的三对象重叠；大图还存在多个可选操作分支，需要进一步搜索操作顺序。枚举方法适合核验小例，随着可交换节点增加，候选数迅速增长，六图因此采用受限候选生成。小例中的最优性由上下界与穷举共同支持，大图的解质量则由实际基准和对应下界评价。

\subsection{候选生成与算法对应}
在操作诱导子图中维护入度和就绪集合。每选择一个操作，先插入尚未出现且为该操作所需的申请，随后发出操作，更新剩余使用次数；只有剩余使用次数为零且全部显式前驱已完成时才插入释放。输出完毕后逐边检查式\eqref{eq:topo}，并检查节点置换及累计量首尾归零。

为了分开不同调度偏好的效果，设置三类主要操作顺序。编号顺序每次选择最小编号的就绪节点，是经过生命期收缩的强基准，而非未经处理的文件行顺序。缓存贪心先选就绪集合中编号最小的 $W$ 个节点，再最小化
\begin{equation}
q_v=n_v-\alpha d_v,\quad
n_v=\sum_{b\in B_v\setminus B_{\rm live}}s_b,\quad
d_v=\sum_{b\in B_v:\,\operatorname{remain}(b)=1}s_b.\label{eq:greedy}
\end{equation}
其中 $d_v$ 是末次使用带来的释放潜力；若还有额外释放前驱未完成，实际释放仍须等待，评分不强制兑现这一潜力。相同评分按新增量、下游路径及编号打破平局。比较 $(W,\alpha)=(32,1),(128,1),(64,2)$，窗口并非从所有就绪节点中寻找评分最好的 $W$ 项。

第三类从编号有序的终端操作出发，沿前驱作深度优先访问，以完成访问顺序得到拓扑序；前驱正序和逆序各形成一个候选。它有机会缩短单个依赖分支的持续时间，但共享缓冲区可能使分支优先与复用优先发生冲突。另保留两种流水评分产生的操作顺序参与峰值比较。最终八种顺序均经过相同生命期收缩，再按真实峰值选优，不用局部评分直接替代目标函数。

令 $Q=\sum_{v\in V_o}|B_v|$，$K=\max_v|B_v|$。图遍历及未排序的深度优先核心为 $O(N+|E|)$；实现中的编号排序另有至多 $O((N+|E|)\log N)$ 开销。窗口贪心每步扫描就绪集合，并计算至多 $W$ 个缓冲区评分，保守界为 $O(N^2\log(W+1)+NWK+|E|+Q)$。空间为 $O(N+|E|+Q)$。实际并未枚举所有拓扑序，因此计算可承受性以限制搜索范围为代价。

\subsection{状态更新与停止条件}
算法的输入是图、窗口 $W$ 和释放权重 $\alpha$，输出是包含全部原节点的序列。状态由已输出节点集合、操作剩余入度、就绪集合、逻辑活跃对象和剩余使用次数组成。操作剩余入度只统计操作间的边；管理依赖在输出申请和释放时单独检查。逆拓扑遍历预先计算下游周期长度，用于评分相同时优先处理后续链条较长的操作。

每次迭代先从就绪集合构造候选窗口，按所选评分确定操作 $v$。随后插入 $v$ 尚缺的缓冲区申请，以及指向 $v$ 而尚未输出的其他申请前驱。全部前驱满足后输出 $v$，将其引用对象的剩余使用次数各减一，并检查因此就绪的释放。最后更新操作后继的入度，将降至零的后继加入就绪集合。表\ref{tab:q1steps}列出这些步骤各自维持的条件。
\begin{table}[H]\centering\small
\caption{操作顺序生成中的状态更新}\label{tab:q1steps}
\begin{tabular}{p{.18\linewidth}p{.36\linewidth}p{.34\linewidth}}\toprule
步骤&更新内容&维持条件\\\midrule
选择操作&按窗口与评分选取就绪节点&全部操作前驱已经输出\\
补齐申请&输出未申请的需求对象&操作发生前对象均活跃\\
执行与计数&输出操作，递减使用计数&计数等于未输出使用者数量\\
检查释放&计数归零且显式前驱完成才释放&后续操作仍需对象保持活跃\\
扩展前沿&后继入度归零后进入就绪集&每个操作恰好输出一次\\\bottomrule
\end{tabular}\end{table}

逻辑活跃集合与剩余使用计数承担不同职责。首次使用对象需要插入申请，已经分配但未来仍要使用的对象持续保留；计数归零时还要检查显式释放前驱。这样即使图中存在额外管理边，释放也不会越过约束。遍历到就绪集合为空时，处理剩余的纯管理节点，并以全节点唯一、全部边正序及最终驻留为零作为完成条件。

例如，当前驻留量为 10，两个就绪操作的 $(n_v,d_v)$ 分别为 $(4,6)$ 和 $(1,2)$。当 $\alpha=1$ 时，两者净评分为 $-2$ 和 $-1$，缓存贪心优先前者；按式\eqref{eq:steppeak}，本步申请峰值却分别为 14 和 11。前者操作后的剩余驻留更小，后者瞬时申请压力更小，局部准则在这两个倾向之间作出取舍。后续哪个分支会带来更多驻留，还取决于尚未执行的使用者，因此最终必须以完整序列的最大值评价。

\subsection{六图峰值与解质量}
表\ref{tab:peaks}同时给出收缩后编号基准、选定峰值及式\eqref{eq:hlower}的下界。比值 $H_1/L_H$ 是相对一个可计算下界的比值，不是已知的最优性误差。
\begin{table}[H]\centering\caption{驻留峰值与同时使用下界（题面容量单位）}\label{tab:peaks}\small
\begin{tabular}{lrrrr}\toprule
计算图 & $H_{\rm id}$ & $H_1$ & $L_H$ & $H_1/L_H$\\\midrule
卷积小图 & 7560 & 6984 & 3144 & 2.22\\
卷积大图 & 14048 & 14048 & 932 & 15.07\\
注意力小图 & 4884 & 4884 & 512 & 9.54\\
注意力大图 & 7972 & 7972 & 512 & 15.57\\
矩阵乘小图 & 9728 & 9728 & 512 & 19.00\\
矩阵乘大图 & 35328 & 35328 & 512 & 69.00\\
\bottomrule\end{tabular}
\end{table}

卷积小图通过深度优先顺序从 7560 降至 6984，减少 576，降幅为 7.62\%。其余五图在本轮候选中保留编号基准，选定峰值分别为 14048、4884、7972、9728 和 35328。矩阵乘大图的同时使用下界为 512，与 35328 的可行峰值之间仍有较大距离。该界主要刻画单次操作，跨操作保留对象则由后续驻留组成进一步说明。

\subsection{驻留轨迹与峰值形成位置}
为观察顺序改变后的过程，图\ref{fig:convtrace}给出卷积小图三种候选的完整驻留轨迹，并放大编号与深度优先方案的峰值附近。横轴为已处理的原始节点位置，纵轴为所有类型容量之和。同一横坐标表示相同的已处理节点数量，两种序列在该位置的具体节点可以不同；因此比较首先围绕各自的峰值及组成展开。
\begin{figure}[H]\centering
\includegraphics[width=.94\linewidth]{../../figures/npu2025-r3/conv_trajectory.pdf}
\caption{卷积小图的驻留轨迹与峰值窗口}\label{fig:convtrace}\end{figure}

编号基准在第 2286 个节点之后达到 7560，此时节点为 1485，活跃对象共 13 个，其中一级缓存占 7488，零级右输入缓存占 72。深度优先方案在第 2278 个节点之后达到 6984，此时节点为 1483，活跃对象共 9 个，两个类型的容量分别为 6912 和 72。两方案峰值处的零级贡献相同，总峰值的 576 差值与一级缓存贡献的差值一致。这一分解把总体收益定位到峰值窗口的类型组成，进一步逐对象追踪才可判断每个对象的申请或释放变化。

局部缓存贪心在 $(W,\alpha)=(32,1)$ 时达到 21762 的峰值，出现在第 1618 个节点之后。该处活跃对象增至 283 个，一级缓存占 5378，而零级右输入缓存占 16384。由此可见，较大的峰值并非仅将原基准的一级缓存高峰继续抬升，而是伴随另一缓存类型的大量逻辑对象同时存续。问题一尚未进行物理分配，16384 表示该顺序要求的逻辑驻留量；后续引入物理容量时，必须通过换入换出处理这种压力。

这与净增量评分的结构相吻合。评分针对当前就绪操作计算首次需求与释放潜力，未显式惩罚某一类型已积累的大量活跃对象。连续选择单步代价较低的操作，可能保留许多还需其他分支使用的中间缓冲区。图中的高峰和对象数量为这一风险提供了具体过程证据，因而候选筛选采用真实峰值，而非累加局部评分。

\begin{table}[H]\centering\small
\caption{各选定序列总峰值发生位置的分类型组成}\label{tab:peaktypes}
\begin{tabular}{lrrrrr}\toprule
计算图 & $L1$ & $UB$ & $L0A$ & $L0B$ & $L0C$\\\midrule
卷积小图 & 6912 & 0 & 0 & 72 & 0\\
卷积大图 & 14040 & 0 & 0 & 8 & 0\\
注意力小图 & 3072 & 1044 & 128 & 128 & 512\\
注意力大图 & 5120 & 2084 & 128 & 128 & 512\\
矩阵乘小图 & 9216 & 0 & 128 & 128 & 256\\
矩阵乘大图 & 34816 & 0 & 128 & 128 & 256\\
\bottomrule\end{tabular}
\end{table}
表\ref{tab:peaktypes}中的每一行取自该图总峰值的同一节点位置，行内相加恰为 $H_1$。卷积大图的 14048 中有 14040 来自一级缓存；矩阵乘两图的一级缓存贡献分别为 9216、34816。注意力两图则同时出现一级缓存和统一缓存的较高占用，小图在总峰值位置的统一缓存需求为 1044，大图为 2084，均高于后续问题二的 1024 容量。这些结果说明，分类型压力对连续分配具有直接意义，单个总峰值不足以决定哪类对象需要换出。

问题一由此给出六份可行原始序列、对应峰值及可计算下界，并揭示编号基准、分支顺序和局部贪心的差异。进入问题二时，保留不同顺序候选用于真实地址分配，使其价值由额外搬运量重新评价。这样既利用了生命期收缩，也保留了不同目标下的选择余地。
